\chapter{Decoding}
[This is only a rough draft! Will need to reread all the papers and then write something fitting!!! with citations]

After extracting the syndrome we need to determine a fitting correction given the extracted syndrome.
This process it called decoding.
(And it is also a problem in classical error correction) 

Given that (in most cases) there exist no 1-to-1 mapping between errors, we need to a way to choose a correction from the set of possible corrections. 

This process can be as simple as a lookup table matching each syndrome with a given correction.
But due to the exponential scaling of the total number of possible syndromes \textcolor{red}{put reference in which I describe the scaling of syndromes here},
it is for larger codes or codes with larger distances more reasonable to calculate the fitting correction on the fly, opposed to precalculating a whole lookup table.

There exist many different approaches to determine a fitting correction given a syndrome.
Choosing a decoder for a given QEC cycle is a process of weighting the benefits and disadvantages. 
In this thesis we will use two different approaches.

First we will introduce the Maximum Likelihood (ML) decoder. 
This decoder will return the correction with the highest likelihood of a successful recovery given the syndrome, hence its name. 
We will go into more detail in the following section \ref{sec:ML_decoding}.

The second decoder we will discuss and use is the Minimum Weight Perfect Matching (MWPM) Decoder.
This decoder return the correction fitting to the most probable single error process.
Further detail of the inner workings will be discussed in section \ref{sec:MWPM_decoding}.



\section{Maximum Likelihood Decoding}\label{sec:ML_decoding}

We will first shortly introduce the general concept of ML Decoding.
Then we will talk about the specific approach we are going to use to decode the syndromes extracted from our circuit with a surface code encoding and Steane type syndrome extraction.

\subsection{General ML Decoder}

[Do an algo box here?]



\subsection{ML Decoder used in this thesis}


\section{Minimum Weight Perfect Matching Decoding}\label{sec:MWPM_decoding}


\section{Pauli Frame Tracking}